% CHTC Researcher Forum 2026: Posters are 4 feet (48 inches) wide and 3 feet (36 inches) tall.
% Hack: Currently set in betterposter.cls
\documentclass[a0paper,fleqn]{betterposter}
\usepackage{bm}
\graphicspath{{figures/}}
\usepackage[hidelinks]{hyperref}

% Uncomment the following commands to customise the format

%% Setting the width of columns
% Left column
\setlength{\leftbarwidth}{0.2\paperwidth}
% Right column
\setlength{\rightbarwidth}{0.2\paperwidth}

%% Setting the column margins
% Horizontal margin
%\setlength{\columnmarginvertical}{0.05\paperheight}
% Vertical margin
%\setlength{\columnmarginhorizontal}{0.05\paperheight}
% Horizontal margin for the main column
%\setlength{\maincolumnmarginvertical}{0.15\paperheight}
% Vertical margin for the main column
%\setlength{\maincolumnmarginhorizontal}{0.15\paperheight}

% Text font
\renewcommand{\fontsizestandard}{\fontsize{32}{46} \selectfont}
% Main column font
% \renewcommand{\fontsizemain}{\fontsize{148.00}{280.00} \selectfont}
% Set baseline skip to 1.4 times font size
% \renewcommand{\fontsizemain}{\fontsize{135}{162} \selectfont}
% \renewcommand{\fontsizemain}{\fontsize{130}{156} \selectfont}
\renewcommand{\fontsizemain}{\fontsize{128}{180} \selectfont}
% \renewcommand{\fontsizemain}{\fontsize{130}{182} \selectfont}
% Title font
% \renewcommand{\fontsizetitle}{\fontsize{101.50}{152.00} \selectfont}
\renewcommand{\fontsizetitle}{\fontsize{80}{112} \selectfont}
% Author font
\renewcommand{\fontsizeauthor}{\fontsize{33.55}{47.3} \selectfont}
\newcommand{\fontsizeinstitution}{\fontsize{20}{25} \selectfont}
% Section font
\renewcommand{\fontsizesection}{\fontsize{61.00}{86.00} \selectfont}
\renewcommand{\fontsizesubsection}{\fontsize{48.00}{57.00} \selectfont}

% Changing font sizes for a specific text segment
% Place the text inside brackets:
% {\fontsize{28}{35} \selectfont Your text goes here}

%% Changing colours
% Background of side columns
%\renewcommand{\columnbackgroundcolor}{black}
% Font of side columns
%\renewcommand{\columnfontcolor}{gray}
% Background of main column
%\renewcommand{\maincolumnbackgroundcolor}{empirical}
%\renewcommand{\maincolumnbackgroundcolor}{theory}
%\renewcommand{\maincolumnbackgroundcolor}{methods}
%\renewcommand{\maincolumnbackgroundcolor}{intervention}
% Font of main column
%\renewcommand{\maincolumnfontcolor}{gray}

% Disable hyphenation
\hyphenpenalty=10000
\exhyphenpenalty=10000

\begin{document}
\betterposter{
 % MAIN COLUMN
 \maincolumn{
  % Main space
  % TODO: Fix alignment with Claude
  % CHTC resources, \\Linux container images, \\and Pixi make \textbf{reproducible GPU workflows} for complex science \textbf{easy}
  \textbf{Reproducible GPU environments} with Pixi, Linux container images, and CHTC resources make complex scientific workflows \textbf{easy to run}
  % CHTC resources, \\Linux container images, \\and Pixi make  for complex science \textbf{easy}
  % \textbf{pyhf} is a \textbf{pure Python} statistical fitting library that uses \textbf{tensors} and \textbf{autograd} to speed up physics analysis at the LHC
 }{
  % Bottom space

\begin{center}
  % QR code
  \qrcode{qr_code.pdf}{smartphoneWhite}{
   \textbf{Take a picture} to get more
   \\information, examples, and this poster\\
   \href{https://github.com/matthewfeickert/poster-chtc-researcher-forum-2026}{github.com/matthewfeickert/poster-chtc-researcher-forum-2026}
  }
\end{center}
}

}{
 % LEFT COLUMN

%  \title{Accelerating\\hardware-accelerated\\science}
 \title{Accelerating\\hardware\\accelerated\\science}
 \vspace{-1em}
%  \textbf{Reproducible hardware-accelerated science}

%  \author{{\href{https://www.matthewfeickert.com/}{\underline{Matthew Feickert}$^{1}$}}, \href{https://dsi.wisc.edu/staff/cranmer-kyle/}{Kyle Cranmer$^{1}$}}
 \author{{\href{https://www.matthewfeickert.com/}{\underline{Matthew Feickert}$^{1}$}}}
 \institution{1 University of Wisconsin--Madison Data Science Institute}
 %
 \section{Hardware-accelerated Science}
 Machine learning is becoming a standard part of the modern scientist's toolkit.
 As a result, the use of \textbf{hardware accelerators} (e.g. GPUs) in scientific software workflows is increasingly common.
 Historically, GPU-enabled software could be difficult to work with and required expertise to deploy.

%  One of the most widely used statistical models in \textbf{high energy physics} for binned measurements and searches

%  \begin{center}
%   \includegraphics[width=\textwidth]{HistFactory_result_examples.png}
%  \end{center}
%  %
%  \vspace{-0.5em}
%  \begin{minipage}{0.29\textwidth}
%   \begin{center}
%    \textbf{Standard Model}
%   \end{center}
%  \end{minipage}%
%  \quad
%  \begin{minipage}{0.36\textwidth}
%   \begin{center}
%    \textbf{Supersymmetry}
%   \end{center}
%  \end{minipage}%
%  \quad
%  \begin{minipage}{0.27\textwidth}
%   \begin{center}
%    \textbf{Exotic Physics}
%   \end{center}
%  \end{minipage}%
%  %
%  \vspace{2em}
%  %

 \vspace{1.0em}
 \textbf{Fully reproducible environments with Pixi}
  \begin{center}
   \href{https://pixi.prefix.dev/}{\includegraphics[width=0.5\textwidth]{tools_pixi.pdf}}
  \end{center}
 \begin{itemize}
  \item Declaratively create fully reproducible locked environments with conda packages (ANY language) and Python packages
  \item Locking: automatic and Rust-fast
  \item Define host-specific expectations of hardware resources (e.g. CUDA version)
  \item Resolve environments for platforms different than solve machine (e.g. solve on Mac laptop, run on Linux machine with CUDA)
 \end{itemize}

 \vspace{1.0em}
 \textbf{Linux container distributions}

 \vspace{1.0em}
 \begin{minipage}{0.50\textwidth}
  \begin{center}
   \href{https://apptainer.org/}{\includegraphics[width=0.8\textwidth]{apptainer_logo.pdf}}
  \end{center}
 \end{minipage}%
 \hfill
 \begin{minipage}{0.50\textwidth}
  \begin{center}
   \href{https://www.docker.com/}{\includegraphics[width=0.8\textwidth]{docker-logo-blue.pdf}}
  \end{center}
 \end{minipage}%
 \vspace{1.0em}

 \begin{itemize}
  % \item Linux container technology (Docker, Apptainer) deploy platform-specific software environments through container images
  \item Docker, Apptainer deploy platform-specific software environments through container images
  \item With Pixi environment, Dockerfile or Apptainer definition files reduce to template
  \item Linux container image for GPU workflows can be automatically built and deployed on code changes
  \item Fast launching with CHTC's \texttt{/staging}
 \end{itemize}

%  \textcolor{blue}{Primary Measurement}:
%  \begin{itemize}
%   \item Multiple disjoint ``channels'' (e.g. event observables) each with multiple bins of data
%   \item Example parameter of interest: strength of physics signal, $\mu$
%  \end{itemize}
%  \textcolor{red}{Auxiliary Measurements}:
%  \begin{itemize}
%   \item Nuisance parameters (e.g. in-situ measurements of background samples)
%   \item Systematic uncertainties (e.g. normalization, shape, luminosity)
%  \end{itemize}

 % \vspace{-1em}
 %
%  \vfill
%  \subsection{Hardware Acceleration}
%  For ML-library tensor backends the computational graph can be transparently placed on hardware accelerators: \textbf{GPUs} and \textbf{TPUs} for order of magnitude speed-up in computation
%  \begin{center}
%   \includegraphics[width=\textwidth]{scaling_hardware.pdf}
%  \end{center}
}{
 % RIGHT COLUMN

 \textbf{HTCondor container universe}
 \vspace{1.0em}
  \begin{center}
   \href{https://htcondor.org/}{\includegraphics[width=0.8\textwidth]{htcondor_logo.pdf}}
  \end{center}
 \vspace{1.0em}
 \begin{itemize}
  \item CHTC's HTCondor deployments support container universe
  \item GPU and CUDA capabilities specified in submit description files
  \item Parallel tasks can scaleout across CHTC GPU nodes
 \end{itemize}

 \subsection{Implementation}
 \begin{center}
  \includegraphics[width=0.2\textwidth]{computational_graph3.pdf}
 \end{center}
 The computational graph of multidimensional array operations for likelihood function of a physics analysis defined through HistFactory

%  %
%  \vspace{0.5em}
%  %
%  \begin{minipage}{0.33\textwidth}
%   \begin{center}
%    \includegraphics[width=\textwidth]{NumPy_logo.pdf}
%   \end{center}
%  \end{minipage}%
%  \quad
%  \begin{minipage}{0.33\textwidth}
%   \begin{center}
%    \includegraphics[width=\textwidth]{TensorFlow_logo.pdf}
%   \end{center}
%  \end{minipage}%
%  \quad
%  \begin{minipage}{0.33\textwidth}
%   \begin{center}
%    \includegraphics[width=0.85\textwidth]{Pytorch_logo.pdf}
%   \end{center}
%  \end{minipage}%
%  %
%  \vspace{0.5em}
%  %

%  Use of $n$-dimensional array (``tensor'') operations through a common API layer around high performance tensor libraries

%  \vspace{-1em}
 \subsection{Applications in DSI Science}
 Data scientists and researchers at the \textbf{University of Wisconsin--Madison Data Science Institute} are using this workflow pattern across:
 \begin{itemize}
  \item High Energy Physics applications of simulation-based Inference
  \item High-Resolution Rapid Refresh (HRRR) weather forecasting
  \item Large-scale simulation of protein motions using molecular dynamics
 \end{itemize}

 \subsection{More Information}
 \begin{itemize}
  \item Matthew Feickert, Ruben Arts, John Kirkham, \emph{Reproducible Machine Learning Workflows for
Scientists with Pixi}, \\\textbf{DOI: \href{https://doi.org/10.25080/nwuf8465}{10.25080/nwuf8465}}
  \item \href{https://github.com/CHTC/templates-GPUs}{\textbf{github.com/CHTC/templates-GPUs}}
  \item \href{https://github.com/CHTC/recipes}{\textbf{github.com/CHTC/recipes}}
 \end{itemize}

%  The full likelihood can be expressed as a \textbf{single JSON document}\\
%  Archive friendly for analysis presentation
%  \vspace{0.5em}
%  \begin{center}
%   \includegraphics[width=0.9\textwidth]{carbon_JSON_spec.png}
%  \end{center}
%  \vspace{-1em}
%  \begin{center}
%   {\fontsizeinstitution\textbf{Example:} 2 binned single channel with 2 samples with 1 parameter of interest and 1 nuisance parameter}
%  \end{center}
%  \begin{center}
%   \includegraphics[width=0.9\textwidth]{carbon_pyhf_CLs.png}
%  \end{center}

 \vspace{1em}
 \vfill
%  \begin{minipage}{0.58\textwidth}
%   \begin{center}
%    \href{http://iris-hep.org/}{\includegraphics[width=\textwidth]{IRIS-HEP_logo}}
%   \end{center}
%  \end{minipage}%
%  \quad
%  \begin{minipage}{0.38\textwidth}
%   \begin{center}
%    \href{https://pypi.org/project/pyhf/}{\includegraphics[width=\textwidth]{pyhf_PyPI.pdf}}
%   \end{center}
%   \vspace{1em}
%   \begin{center}
%    \href{https://doi.org/10.5281/zenodo.1169739}{\includegraphics[width=\textwidth]{zenodo_doi.pdf}}
%   \end{center}
%   \vspace{3em}
%  \end{minipage}%
\begin{center}
  \href{https://dsi.wisc.edu/}{\includegraphics[width=\textwidth]{DSI_logo.png}}
\end{center}
}
\end{document}
